% ============================================================
%  يُترجَم بـ XeLaTeX  (في Overleaf: Menu → Compiler → XeLaTeX)
%  الخط: Amiri (متوفر في Overleaf، أو حمّله من github.com/aliftype/amiri)
% ============================================================
\documentclass[12pt,a4paper]{article}
\usepackage[a4paper,margin=1.7cm,top=2.3cm,bottom=2.9cm,headheight=15pt,footskip=1.5cm]{geometry}
\usepackage{amsmath,amssymb}
\usepackage{icomma}
\usepackage[english,bidi=default,layout=lists]{babel}
\babelprovide[import,main]{arabic}
\babelfont{rm}[Scale=1.08]{Amiri}
\babelfont{sf}[Scale=1.08]{Amiri}
\usepackage[most]{tcolorbox}
\usepackage{tikz}
\usetikzlibrary{arrows.meta,calc,shadings,patterns}
\usepackage{array}
\usepackage{colortbl}
\usepackage{fancyhdr}
\usepackage{titlesec}
\usepackage{enumitem}
\usepackage{collcell}
\usepackage{needspace}
\usepackage{etoolbox}
\newcolumntype{C}{>{\collectcell\ar}c<{\endcollectcell}}
\pretocmd{\section}{\Needspace{15\baselineskip}}{}{}
\pretocmd{\subsection}{\Needspace{11\baselineskip}}{}{}

\definecolor{cDef}{HTML}{1F4E9C}
\definecolor{cThm}{HTML}{B3261E}
\definecolor{cRem}{HTML}{2E7D32}
\definecolor{cEx}{HTML}{C76B00}
\definecolor{cMeth}{HTML}{6A1B9A}
\definecolor{cSol}{HTML}{00695C}
\definecolor{cSec}{HTML}{1F3A68}

\titleformat{\section}{\Large\bfseries\color{cSec}}{\thesection.}{0.6em}{}[{\color{cSec!60}\titlerule[1pt]}]
\titleformat{\subsection}{\large\bfseries\color{cThm!85!black}}{\thesubsection.}{0.5em}{}
\titlespacing*{\section}{0pt}{14pt}{8pt}
\titlespacing*{\subsection}{0pt}{10pt}{5pt}
\setlength{\parindent}{0pt}
\setlength{\parskip}{4pt}
\setlist{itemsep=2pt,topsep=3pt}

\newtcolorbox{arbox}[3][]{enhanced,
  colback=#2!5,colframe=#2!80!black,arc=2mm,boxrule=0.8pt,
  left=6pt,right=6pt,top=8pt,bottom=5pt,
  fonttitle=\bfseries,coltitle=white,
  title={\foreignlanguage{arabic}{#3}},
  attach boxed title to top right={xshift=-5mm,yshift=-3mm},
  boxed title style={colback=#2!80!black,arc=1.5mm,boxrule=0pt},
  before upper={\selectlanguage{arabic}},
  before skip=8pt,after skip=8pt,breakable,#1}
\newenvironment{BOX}[2]{\begin{otherlanguage}{english}\begin{arbox}{#1}{#2}}{\end{arbox}\end{otherlanguage}}
\newenvironment{tarif}[1][تعريف]{\begin{BOX}{cDef}{#1}}{\end{BOX}}
\newenvironment{khassia}[1][خاصية]{\begin{BOX}{cThm}{#1}}{\end{BOX}}
\newenvironment{molahada}[1][ملاحظة]{\begin{BOX}{cRem}{#1}}{\end{BOX}}
\newenvironment{mithal}[1][مثال]{\begin{BOX}{cEx}{#1}}{\end{BOX}}
\newenvironment{tariqa}[1][طريقة]{\begin{BOX}{cMeth}{#1}}{\end{BOX}}

\newcounter{exo}
\newenvironment{tamrin}[1]{\stepcounter{exo}\gdef\curt{#1}\begin{BOX}{cSec}{#1}}{\end{BOX}}
\newcommand{\curt}{}
\newenvironment{hal}{\Needspace{12\baselineskip}\begin{otherlanguage}{english}\begin{arbox}[breakable]{cSol}{حل \curt}}{\end{arbox}\end{otherlanguage}}

\newenvironment{LTR}{\begin{otherlanguage}{english}}{\end{otherlanguage}}
\newcommand{\ar}[1]{\foreignlanguage{arabic}{#1}}

\renewcommand{\dots}{\mathinner{\ldotp\ldotp\ldotp}}
\renewcommand{\ldots}{\mathinner{\ldotp\ldotp\ldotp}}
\renewcommand{\cdots}{\mathinner{\cdotp\cdotp\cdotp}}
\newcommand{\N}{\mathbb{N}}
\newcommand{\Z}{\mathbb{Z}}
\newcommand{\D}{\mathbb{D}}
\newcommand{\Q}{\mathbb{Q}}
\newcommand{\R}{\mathbb{R}}
\newcommand{\pgcd}{\operatorname{PGCD}}
\newcommand{\ppcm}{\operatorname{PPCM}}

% \NLseg{a}{b}{[}{]} ; \NLright{a}{[} ; \NLleft{b}{]}  (interval on a number line)
\newcommand{\NLseg}[4]{\begin{LTR}\begin{tikzpicture}[baseline=-0.6ex]
 \draw[-{Stealth[length=2mm]},gray!70](0,0)--(3.4,0);
 \draw[line width=2pt,cThm!80!black](0.7,0)--(2.7,0);
 \node[cThm!80!black,font=\LARGE] at (0.7,0){$#3$};
 \node[cThm!80!black,font=\LARGE] at (2.7,0){$#4$};
 \node[below=3pt,font=\footnotesize] at (0.7,0){$#1$};
 \node[below=3pt,font=\footnotesize] at (2.7,0){$#2$};
\end{tikzpicture}\end{LTR}}
\newcommand{\NLright}[2]{\begin{LTR}\begin{tikzpicture}[baseline=-0.6ex]
 \draw[-{Stealth[length=2mm]},gray!70](0,0)--(3.4,0);
 \draw[line width=2pt,cThm!80!black](1.3,0)--(3.3,0);
 \node[cThm!80!black,font=\LARGE] at (1.3,0){$#2$};
 \node[below=3pt,font=\footnotesize] at (1.3,0){$#1$};
\end{tikzpicture}\end{LTR}}
\newcommand{\NLleft}[2]{\begin{LTR}\begin{tikzpicture}[baseline=-0.6ex]
 \draw[-{Stealth[length=2mm]},gray!70](0,0)--(3.4,0);
 \draw[line width=2pt,cThm!80!black](0.05,0)--(2.1,0);
 \node[cThm!80!black,font=\LARGE] at (2.1,0){$#2$};
 \node[below=3pt,font=\footnotesize] at (2.1,0){$#1$};
\end{tikzpicture}\end{LTR}}

% ---------- social-media icons (pure TikZ) ----------
\definecolor{cFB}{HTML}{1877F2}
\definecolor{cTG}{HTML}{229ED9}
\definecolor{cIG1}{HTML}{FEDA75}
\definecolor{cIG2}{HTML}{D62976}
\definecolor{cIG3}{HTML}{4F5BD5}
\newcommand{\fbicon}{\begin{LTR}\begin{tikzpicture}[baseline=-0.75ex,x=1em,y=1em]
 \fill[cFB] (0,0) circle (0.55);
 \node[white,inner sep=0,font=\usefont{TU}{lmss}{bx}{n}\small] at (0.04,-0.02) {f};
\end{tikzpicture}\end{LTR}}
\newcommand{\igicon}{\begin{LTR}\begin{tikzpicture}[baseline=-0.75ex,x=1em,y=1em]
 \shade[lower left=cIG1,lower right=cIG2,upper left=cIG2,upper right=cIG3,rounded corners=0.28em] (-0.55,-0.55) rectangle (0.55,0.55);
 \draw[white,line width=0.09em,rounded corners=0.18em] (-0.36,-0.36) rectangle (0.36,0.36);
 \draw[white,line width=0.09em] (0,0) circle (0.17);
 \fill[white] (0.22,0.22) circle (0.05);
\end{tikzpicture}\end{LTR}}
\newcommand{\tgicon}{\begin{LTR}\begin{tikzpicture}[baseline=-0.75ex,x=1em,y=1em]
 \fill[cTG] (0,0) circle (0.55);
 \fill[white] (-0.33,0.0) -- (0.32,0.27) -- (0.2,-0.3) -- (0.0,-0.15) -- (-0.08,-0.28) -- (-0.08,-0.08) -- cycle;
 \fill[cTG!70!black] (-0.08,-0.08) -- (0.25,0.2) -- (0.0,-0.15) -- (-0.08,-0.28) -- cycle;
\end{tikzpicture}\end{LTR}}
% page number: black box, white two-digit number, at the right of the social row
\newcommand{\PageBox}{\begin{LTR}\tikz[baseline=(p.base)]\node[fill=black,text=white,inner sep=2.5pt,minimum width=1.6em,minimum height=1.3em,font=\usefont{TU}{lmss}{bx}{n}\small](p){\ifnum\value{page}<10 0\fi\arabic{page}};\end{LTR}}
% visual left-to-right order: Telegram | Instagram | Facebook | page box
\newcommand{\SocialRow}{\begin{otherlanguage}{english}%
 \ar{\textbf{الأستاذ سوايسية هشام للرياضيات}}\ \tgicon\hfill
 \textbf{prof\_hicham\_math}\ \igicon\hfill
 \ar{\textbf{الأستاذ سوايسية هشام للرياضيات}}\ \fbicon\hspace{2.2em}\PageBox\end{otherlanguage}}

\pagestyle{fancy}
\fancyhf{}
\renewcommand{\headrulewidth}{0.6pt}
\renewcommand{\footrulewidth}{0.4pt}
% \SetHeader{right head: document title}{left head: level -- stream}
\newcommand{\SetHeader}[2]{%
 \fancyhead[C]{\parbox{\textwidth}{\ar{#1\hfill #2}}}%
 \fancyfoot[C]{\parbox{\textwidth}{\small\SocialRow}}}
\usepackage{tkz-tab}
\usepackage{multicol}
% variation / sign table (LTR)
\newenvironment{TAB}{\par\begin{center}\begin{LTR}\begin{tikzpicture}[scale=0.95,every node/.style={font=\small}]}{\end{tikzpicture}\end{LTR}\end{center}}
% grid + axes, clipped: \begin{GR}{xmin}{xmax}{ymin}{ymax}{unit cm}
\newenvironment{GR}[5]{\par\begin{center}\begin{LTR}\begin{tikzpicture}[x=#5cm,y=#5cm,font=\footnotesize]
 \draw[gray!25,very thin] (#1,#3) grid (#2,#4);
 \draw[-{Stealth[length=2mm]}] (#1,0)--(#2+0.3,0) node[below left]{$x$};
 \draw[-{Stealth[length=2mm]}] (0,#3)--(0,#4+0.3) node[below left]{$y$};
 \node[below left] at (0,0){$O$};
 \begin{scope}\clip (#1,#3) rectangle (#2,#4);}
{\end{scope}\end{tikzpicture}\end{LTR}\end{center}}
\newcommand{\Ticks}[2]{\foreach \t in {#1} \node[below,font=\scriptsize] at (\t,0){$\t$}; \foreach \t in {#2} \node[left,font=\scriptsize] at (0,\t){$\t$};}
\newcommand{\bx}[1]{\ensuremath{\colorbox{cDef!12}{$\displaystyle #1$}}}
\newcommand{\TwoCol}[4]{\par\begin{LTR}\noindent\begin{minipage}[t]{#1\textwidth}\selectlanguage{arabic}#2\end{minipage}\hfill\begin{minipage}[t]{#3\textwidth}\selectlanguage{arabic}#4\end{minipage}\end{LTR}\par}
\newcommand{\RL}[4]{\TwoCol{#3}{#4}{#1}{#2}}  % \RL{w1}{right content}{w2}{left content}
\newcommand{\ok}{$\textcolor{cRem}{\checkmark}$}
\newcommand{\limn}{\displaystyle\lim_{n\to+\infty}}
% sub-questions, each on its own line
\newenvironment{sq}{\begin{itemize}[leftmargin=2.4em,itemsep=3pt]}{\end{itemize}}
\newcommand{\qa}{\item[(أ)]}\newcommand{\qb}{\item[(ب)]}\newcommand{\qc}{\item[(جـ)]}\newcommand{\qd}{\item[(د)]}\newcommand{\qe}{\item[(هـ)]}
\newcommand{\pa}{\par\smallskip\textbf{(أ)} }\newcommand{\pb}{\par\smallskip\textbf{(ب)} }\newcommand{\pc}{\par\smallskip\textbf{(جـ)} }\newcommand{\pd}{\par\smallskip\textbf{(د)} }\newcommand{\pe}{\par\smallskip\textbf{(هـ)} }
\newcommand{\QQ}[1]{\par\smallskip\textbf{#1.} }
\newenvironment{tmr}{\stepcounter{exo}\begin{BOX}{cSec}{تمرين \theexo}}{\end{BOX}}
\newcommand{\ABC}[3]{\\[2pt]\begin{tabular}{p{0.28\textwidth}p{0.28\textwidth}p{0.28\textwidth}}(جـ) \ #3 & (ب) \ #2 & (أ) \ #1\end{tabular}}
\newtcbox{\bxtbox}{on line,colback=cDef!12,colframe=cDef!12,boxrule=0pt,arc=1pt,boxsep=0pt,left=2pt,right=2pt,top=1pt,bottom=1pt}
\newcommand{\bxt}[1]{\foreignlanguage{english}{\bxtbox{\ar{#1}}}}
\SetHeader{\small\bfseries\color{cSec}تمارين الدالة الأسية}{\small\color{cSec}السنة الثالثة ثانوي -- تسيير واقتصاد}
\newenvironment{FIG}{\par\begin{center}\begin{LTR}}{\end{LTR}\end{center}}
\renewcommand{\bx}[1]{\ensuremath{\colorbox{cDef!12}{$\displaystyle\color{cThm}#1$}}}
\newcommand{\LW}[1]{{\color{cDef}#1}}
\newcommand{\Car}[1]{\left\{\begin{array}{l}#1\end{array}\right.}
\newcommand{\Lm}[1]{\displaystyle\lim_{x\to#1}}
\newcommand{\LD}[3]{\par لدينا: $\Lm{#2}#1=\bx{#3}$ لأنّ: \quad}
\newcommand{\E}{\par\smallskip}
\relpenalty=10000\binoppenalty=10000\tolerance=9000\emergencystretch=0.3\linewidth
\begin{document}
\begin{LTR}
\begin{tcolorbox}[enhanced,colback=cSec!10,colframe=cSec!40,arc=3mm,halign=center,boxrule=0pt,top=12pt,bottom=12pt,before upper={\selectlanguage{arabic}}]
{\color{cSec}\fontsize{28}{34}\selectfont\bfseries تمارين الدالة الأسية}\\[8pt]
{\color{cSec!80!black}\large تمارين من الاختبارات + مسائل البكالوريا (2017 -- 2026) + اختيار من متعدد  -- بدون حلول (للتدرّب)}\\[5pt]
{\color{cSec!80!black}\large السنة الثالثة ثانوي -- شعبة تسيير واقتصاد}
\end{tcolorbox}
\end{LTR}

% =====================================================
\section{تمارين من الاختبارات (البكالوريا التجريبية)}

\begin{tamrin}{تمرين من بكالوريا تجريبية رقم 1}
$P(x)$ كثير الحدود المعرف على $\R$ بـ: $P(x)=2x^3-6x-4$.
\QQ{1}
\begin{sq}
\qa تحقق أنه من أجل كل عدد حقيقي $x$: $P(x)=(2x-4)(x+1)^2$.
\qb حل في $\R$ المعادلة $P(x)=0$.
\end{sq}
\QQ{2} (بكالوريا 2026) اعتمادًا على نتائج السؤال السابق:
\begin{sq}
\qa حل في $\R$ المعادلة: $2e^{3x}-6e^x-4=0$.
\qb حل في $]0;+\infty[$ المعادلة: $2(\ln x)^3-6\ln x-4=0$.
\qc حل في $\R$ المتراجحة: $2e^{3x}-6e^x-4\leqslant0$.
\end{sq}
\end{tamrin}

\begin{tamrin}{تمرين من بكالوريا تجريبية رقم 2}
نعتبر الدالة $f$ المعرفة على $\R$ بـ: $f(x)=x-\dfrac{e^x-1}{e^x+1}$ ، و $(C_f)$ تمثيلها البياني في معلم متعامد ومتجانس.
\QQ{1} (بكالوريا 2026) أحسب نهايات الدالة $f$ عند أطراف مجموعة تعريفها.
\QQ{2}
\begin{sq}
\qa بيّن أنه من أجل كل عدد حقيقي $x$: $f'(x)=\dfrac{e^{2x}+1}{(e^x+1)^2}$.
\qb استنتج اتجاه تغير الدالة $f$ ، ثم شكّل جدول تغيراتها.
\end{sq}
\QQ{3}
\begin{sq}
\qa أثبت أنه من أجل كل عدد حقيقي $x$: $f(x)=x-1+\dfrac2{e^x+1}$.
\qb بيّن أن المستقيم $(\Delta)$ ذا المعادلة $y=x-1$ مقارب مائل لـ $(C_f)$ بجوار $+\infty$ ، ثم ادرس وضعية $(C_f)$ بالنسبة إلى $(\Delta)$.
\end{sq}
\QQ{4}
\begin{sq}
\qa أثبت أنه من أجل كل عدد حقيقي $x$: $f(x)=x+1-\dfrac{2e^x}{e^x+1}$.
\qb بيّن أن المستقيم $(\Delta')$ ذا المعادلة $y=x+1$ مقارب مائل لـ $(C_f)$ بجوار $-\infty$ ، ثم ادرس الوضع النسبي لـ $(C_f)$ و $(\Delta')$.
\end{sq}
\QQ{5} (بكالوريا 2025) أحسب $f(0)$ ، ثم ارسم كلًا من $(\Delta)$ و $(\Delta')$ ثم $(C_f)$.
\QQ{6} (بكالوريا 2024) $(\star)$
\begin{sq}
\qa تحقق أن الدالة $F$ المعرفة على $\R$ بـ: $F(x)=\dfrac{x^2}2+x-2\ln(e^x+1)$ دالة أصلية للدالة $f$.
\qb أحسب القيمة المتوسطة للدالة $f$ على المجال $[0;2]$.
\end{sq}
\end{tamrin}

\begin{tamrin}{تمرين من بكالوريا تجريبية رقم 3}
$f$ الدالة المعرفة على $\R$ بـ: $f(x)=\dfrac{4e^{-x}}{e^{-x}+1}-3$ ، و $(C_f)$ تمثيلها البياني.
\QQ{1}
\begin{sq}
\qa بيّن أنه من أجل كل عدد حقيقي $x$: $f(x)=\dfrac4{e^x+1}-3$.
\qb أحسب نهاية الدالة $f$ عند $+\infty$ وعند $-\infty$ ، ثم فسّر النتيجتين هندسيًا.
\end{sq}
\QQ{2} ادرس اتجاه تغير الدالة $f$ ، ثم شكّل جدول تغيراتها.
\QQ{3}
\begin{sq}
\qa جد فاصلة نقطة تقاطع $(C_f)$ مع محور الفواصل.
\qb أكتب معادلة المماس $(T)$ للمنحنى $(C_f)$ في النقطة $\Omega(0;-1)$.
\qc بيّن أنه من أجل كل عدد حقيقي $x$: $f(x)+f(-x)=-2$ ، ثم استنتج أن $(C_f)$ يقبل مركز تناظر.
\qd ارسم $(T)$ و $(C_f)$.
\end{sq}
\QQ{4} $(\star)$ (بكالوريا 2026) أحسب مساحة الحيز المستوي المحدد بالمنحنى $(C_f)$ والمستقيمات: $x=-\ln3$ ، $x=0$ و $y=0$.
\QQ{5} $h$ الدالة المعرفة على $\R$ بـ: $h(x)=f(|x|)$ ، و $(C_h)$ منحناها.
\begin{sq}
\qa بيّن أن $h$ دالة زوجية.
\qb اعتمادًا على $(C_f)$ اشرح كيف يتم رسم $(C_h)$ ، ثم ارسمه.
\end{sq}
\end{tamrin}

\begin{tamrin}{تمرين من بكالوريا تجريبية رقم 4}
\textbf{I)} لتكن الدالة $g$ المعرفة على $\R$ بـ: $g(x)=e^x-x+2$.
\QQ{1} أحسب $\Lm{-\infty}g(x)$ و $\Lm{+\infty}g(x)$.
\QQ{2} ادرس اتجاه تغير الدالة $g$ وشكّل جدول تغيراتها.
\QQ{3} (بكالوريا 2026) استنتج حسب قيم العدد الحقيقي $x$ إشارة $g(x)$.\par
\textbf{II)} نعتبر الدالة $f$ المعرفة على $\R$ بـ: $f(x)=x+(x-1)e^{-x}$ ، و $(C_f)$ منحناها البياني.
\QQ{1} أحسب $\Lm{-\infty}f(x)$ و $\Lm{+\infty}f(x)$.
\QQ{2}
\begin{sq}
\qa بيّن أنه من أجل كل $x$ من $\R$: $f'(x)=e^{-x}g(x)$.
\qb استنتج اتجاه تغير الدالة $f$ وشكّل جدول تغيراتها.
\end{sq}
\QQ{3}
\begin{sq}
\qa بيّن أن المستقيم $(\Delta)$ ذا المعادلة $y=x$ مستقيم مقارب مائل لـ $(C_f)$ عند $+\infty$.
\qb ادرس وضعية $(C_f)$ بالنسبة إلى $(\Delta)$.
\end{sq}
\QQ{4} بيّن أن المعادلة $f(x)=0$ تقبل حلًا وحيدًا $\alpha$ حيث $0{,}3<\alpha<0{,}5$ ، ثم فسّر النتيجة هندسيًا.
\QQ{5} أنشئ $(\Delta)$ و $(C_f)$.
\QQ{6} $(\star)$
\begin{sq}
\qa بيّن أنه من أجل كل $x$ من $\R$: $f(x)-x=1+e^{-x}-f'(x)$ ، ثم استنتج دالة أصلية للدالة $x\mapsto f(x)-x$ على $\R$.
\qb أحسب بدلالة $\alpha$ مساحة الحيز المستوي المحدد بالمنحنى $(C_f)$ والمستقيم $(\Delta)$ والمستقيمين $x=\alpha$ و $x=1$.
\end{sq}
\end{tamrin}

% =====================================================
\section{مسائل البكالوريا -- شعبة تسيير واقتصاد}

\begin{tamrin}{تمرين من بكالوريا 2026 -- الموضوع الثاني}
$f$ الدالة المعرفة على $\R$ بـ: $f(x)=x^2+2x+2xe^x$ ، و $(c_f)$ تمثيلها البياني.
\QQ{1} أحسب $\Lm{-\infty}f(x)$ وبيّن أن: $\Lm{+\infty}f(x)=+\infty$.
\QQ{2}
\begin{sq}
\qa بيّن أنه من أجل كل $x$ من $\R$: $f'(x)=2(x+1)(e^x+1)$.
\qb استنتج اتجاه تغير الدالة $f$ ثم شكّل جدول تغيراتها.
\end{sq}
\QQ{3} عيّن معادلة لـ $(T)$ مماس المنحنى $(c_f)$ عند النقطة ذات الفاصلة $0$.
\QQ{4} $g$ الدالة المعرفة على $\R$ بـ: $g(x)=x^2$ و $(c_g)$ تمثيلها البياني. ادرس الوضع النسبي لـ $(c_f)$ و $(c_g)$.
\QQ{5} أحسب $f(-3)$ ، $f(-2)$ ، $f(1)$ ثم ارسم $(T)$ ، $(c_g)$ و $(c_f)$.
\QQ{6} $(\star)$
\begin{sq}
\qa تحقق أن الدالة $H:x\mapsto(x-1)e^x$ أصلية للدالة $h:x\mapsto xe^x$ على $\R$.
\qb استنتج حساب $\mathcal A$ مساحة الحيز المحدد بالمنحنيين $(c_f)$ و $(c_g)$ والمستقيمين $x=0$ و $x=1$.
\end{sq}
\end{tamrin}

\begin{tamrin}{تمرين من بكالوريا 2025 -- الموضوع الثاني}
\textbf{I)} $g$ الدالة المعرفة على $\R$ بـ: $g(x)=1-(x+1)e^x$.
\QQ{1} أحسب $\Lm{+\infty}g(x)$ وبيّن أن: $\Lm{-\infty}g(x)=1$.
\QQ{2}
\begin{sq}
\qa تحقق أنه من أجل كل عدد حقيقي $x$: $g'(x)=-(x+2)e^x$.
\qb استنتج اتجاه تغير الدالة $g$ ثم شكّل جدول تغيراتها.
\end{sq}
\QQ{3} أحسب $g(0)$ ثم استنتج حسب قيم $x$ إشارة $g(x)$.\par
\textbf{II)} $f$ الدالة المعرفة على $\R$ بـ: $f(x)=x+2-xe^x$ ، و $(c_f)$ تمثيلها البياني (وحدة الطول $2\,cm$).
\QQ{1}
\begin{sq}
\qa أحسب $\Lm{+\infty}f(x)$ وبيّن أن: $\Lm{-\infty}f(x)=-\infty$.
\qb بيّن أن المستقيم $(\Delta)$ ذا المعادلة $y=x+2$ مقارب مائل للمنحنى $(c_f)$ عند $-\infty$.
\qc ادرس الوضع النسبي للمنحنى $(c_f)$ والمستقيم $(\Delta)$.
\end{sq}
\QQ{2}
\begin{sq}
\qa بيّن أنه من أجل كل عدد حقيقي $x$: $f'(x)=g(x)$.
\qb استنتج اتجاه تغير الدالة $f$ ثم شكّل جدول تغيراتها.
\end{sq}
\QQ{3} بيّن أن المنحنى $(c_f)$ يقبل مماسًا $(T)$ موازيًا لـ $(\Delta)$ ، يطلب تعيين معادلة له.
\QQ{4}
\begin{sq}
\qa أحسب $f(1)$ ، $f(-2)$ ثم ارسم كلًا من $(\Delta)$ ، $(T)$ و $(c_f)$.
\qb ناقش بيانيًا وحسب قيم العدد الحقيقي $m$ عدد حلول المعادلة: $f(x)=x+m$.
\end{sq}
\QQ{5} $(\star)$ $h$ الدالة المعرفة على $\R$ بـ: $h(x)=(1-x)e^x$. أحسب $h'(x)$ ثم استنتج بالسنتيمتر المربع مساحة الحيز المحدد بالمنحنى $(c_f)$ والمستقيمات $y=x+2$ ، $x=-1$ و $x=0$.
\end{tamrin}

\begin{tamrin}{تمرين من بكالوريا 2024 -- الموضوع الأول}
\textbf{I)} $g$ الدالة المعرفة على $\R$ بـ: $g(x)=1-(x+1)e^{-x}$.
\QQ{1} ادرس اتجاه تغير الدالة $g$.
\QQ{2} أحسب $g(0)$ ثم استنتج أنه: من أجل كل عدد حقيقي $x$: $g(x)\geqslant0$.\par
\textbf{II)} $f$ الدالة المعرفة على $\R$ بـ: $f(x)=x+(x+2)e^{-x}$ ، و $(C_f)$ تمثيلها البياني.
\QQ{1} أحسب $\Lm{-\infty}f(x)$ ثم بيّن أن: $\Lm{+\infty}f(x)=+\infty$.
\QQ{2}
\begin{sq}
\qa بيّن أنه من أجل كل عدد حقيقي $x$: $f'(x)=g(x)$.
\qb استنتج اتجاه تغير الدالة $f$ ثم شكّل جدول تغيراتها.
\end{sq}
\QQ{3} اكتب معادلة لـ $(T)$ مماس المنحنى $(C_f)$ في النقطة ذات الفاصلة $0$.
\QQ{4}
\begin{sq}
\qa بيّن أن المستقيم $(\Delta)$ ذا المعادلة $y=x$ مقارب مائل للمنحنى $(C_f)$ عند $+\infty$.
\qb ادرس الوضع النسبي للمنحنى $(C_f)$ والمستقيم $(\Delta)$.
\end{sq}
\QQ{5}
\begin{sq}
\qa بيّن أن $(C_f)$ يقطع حامل محور الفواصل في نقطة وحيدة فاصلتها $\alpha$ حيث $-1{,}69<\alpha<-1{,}68$.
\qb ارسم $(T)$ و $(\Delta)$ ثم $(C_f)$.
\end{sq}
\QQ{6} $(\star)$
\begin{sq}
\qa بيّن أن الدالة $H:x\mapsto-(x+3)e^{-x}$ هي دالة أصلية للدالة $h:x\mapsto(x+2)e^{-x}$ على $\R$.
\qb أحسب $\mathcal A$ مساحة الحيز المستوي المحدد بالمنحنى $(C_f)$ والمستقيمات: $y=x$ ، $x=0$ و $x=2$.
\end{sq}
\end{tamrin}

\begin{tamrin}{تمرين من بكالوريا 2023 -- الموضوع الأول}
$f$ الدالة المعرفة على المجال $[0;+\infty[$ بـ: $f(x)=x+1-\dfrac3{e^x+1}$ ، و $(C_f)$ تمثيلها البياني (وحدة الطول $2\,cm$).
\QQ{1} بيّن أن: $\Lm{+\infty}f(x)=+\infty$.
\QQ{2}
\begin{sq}
\qa بيّن أن المستقيم $(\Delta)$ ذا المعادلة $y=x+1$ مقارب مائل لـ $(C_f)$ عند $+\infty$.
\qb ادرس وضعية $(C_f)$ بالنسبة إلى $(\Delta)$.
\end{sq}
\QQ{3}
\begin{sq}
\qa بيّن أنه من أجل كل $x$ من $[0;+\infty[$: $f'(x)=1+\dfrac{3e^x}{(e^x+1)^2}$.
\qb استنتج أن الدالة $f$ متزايدة تمامًا على $[0;+\infty[$ ثم شكّل جدول تغيراتها.
\end{sq}
\QQ{4} بيّن أن المعادلة $f(x)=0$ تقبل حلًا وحيدًا $\alpha$ حيث $0{,}28<\alpha<0{,}29$.
\QQ{5} ارسم $(\Delta)$ و $(C_f)$.
\QQ{6} $(\star)$ $F$ الدالة المعرفة على $[0;+\infty[$ بـ: $F(x)=3x-3\ln(e^x+1)$.
\begin{sq}
\qa تحقق أن $F$ أصلية للدالة $x\mapsto\dfrac3{e^x+1}$ على $[0;+\infty[$.
\qb أحسب بالسنتيمتر المربع مساحة الحيز المحدد بالمنحنى $(C_f)$ والمستقيم $(\Delta)$ والمستقيمين $x=0$ و $x=\ln2$.
\end{sq}
\end{tamrin}

\begin{tamrin}{تمرين من بكالوريا 2019 -- الموضوع الثاني}
\textbf{I)} $g$ الدالة العددية المعرفة على المجال $]-\infty;0]$ كما يلي: $g(x)=2x+6-e^{2x+1}$.
\QQ{1}
\begin{sq}
\qa أحسب $\Lm{-\infty}g(x)$.
\qb ادرس اتجاه تغير $g$ على المجال $]-\infty;0]$ ثم شكّل جدول تغيراتها.
\end{sq}
\QQ{2}
\begin{sq}
\qa بيّن أن المعادلة $g(x)=0$ تقبل حلًا وحيدًا $\alpha$ حيث $-3<\alpha<-2{,}9$.
\qb استنتج إشارة $g(x)$ على المجال $]-\infty;0]$.
\end{sq}
\textbf{II)} $f$ الدالة المعرفة على المجال $]-\infty;0]$ كما يلي: $f(x)=-2x^2-12x+e^{2x+1}$ ، و $(C_f)$ تمثيلها البياني (الوحدة $1\,cm$ على محور الفواصل و $0{,}5\,cm$ على محور التراتيب).
\QQ{1} أثبت أنه من أجل كل $x$ من $]-\infty;0]$: $f'(x)=-2g(x)$.
\QQ{2} استنتج اتجاه تغير الدالة $f$ على المجال $]-\infty;0]$.
\QQ{3} أحسب $\Lm{-\infty}f(x)$ ثم شكّل جدول التغيرات للدالة $f$.
\QQ{4} بيّن أن: $f(\alpha)=-2\alpha(\alpha+5)+6$ وأعط حصرًا للعدد $f(\alpha)$ ، ثم ارسم $(C_f)$ على المجال $[-4;0]$.
\QQ{5} $(\star)$ أحسب بدلالة $\alpha$ التكامل $\dfrac12\displaystyle\int_\alpha^0f(x)\,dx$ ثم فسّر النتيجة بيانيًا.
\end{tamrin}

\begin{tamrin}{تمرين من بكالوريا 2018 -- الموضوع الثاني}
\textbf{I)} نعتبر الدالة العددية $g$ المعرفة على $[0;+\infty[$ بـ: $g(x)=1+(1-x)e^{-x+1}$.\par
ادرس اتجاه تغير الدالة $g$ ثم بيّن أنه من أجل كل $x$ من $[0;+\infty[$: $g(x)>0$ (لا يطلب حساب النهايات).\par
\textbf{II)} نعتبر الدالة العددية $f$ المعرفة على $[0;+\infty[$ بـ: $f(x)=x+xe^{-x+1}$ ، و $(C_f)$ تمثيلها البياني.
\QQ{1}
\begin{sq}
\qa أحسب $\Lm{+\infty}f(x)$ ، ثم بيّن أن المستقيم $(\Delta)$ ذا المعادلة $y=x$ مقارب للمنحنى $(C_f)$.
\qb ادرس وضعية المنحنى $(C_f)$ بالنسبة إلى المستقيم $(\Delta)$.
\end{sq}
\QQ{2} بيّن أنه من أجل كل $x$ من $[0;+\infty[$: $f'(x)=g(x)$ ، ثم شكّل جدول التغيرات للدالة $f$.
\QQ{3} بيّن أن المعادلة $f(x)=4$ تقبل حلًا وحيدًا $\alpha$ حيث: $3{,}75<\alpha<3{,}77$.
\QQ{4} اكتب معادلة المماس $(T)$ للمنحنى $(C_f)$ عند النقطة ذات الفاصلة $1$ ، ثم ارسم $(T)$ ، $(\Delta)$ و $(C_f)$.
\QQ{5} $(\star)$ نعتبر الدالة $F$ المعرفة على $[0;+\infty[$ بـ: $F(x)=\dfrac12x^2-(x+1)e^{-x+1}$.
\begin{sq}
\qa بيّن أن الدالة $F$ هي دالة أصلية للدالة $f$ على المجال $[0;+\infty[$.
\qb أوجد القيمة المضبوطة للعدد $\displaystyle\int_1^4f(x)\,dx$ ، ثم أعط تفسيرًا هندسيًا لهذا العدد.
\end{sq}
\QQ{6} تنمذج الكلفة الهامشية $C_m$ لإنتاج كمية $q$ (مقدرة بآلاف الوحدات) حيث $0\leqslant q\leqslant7$ بالدالة $f$ ، أي: $C_m(q)=f(q)$ (الكلفة الهامشية مقدرة بملايين الدنانير).
\begin{sq}
\qa ما هي كمية المنتوج التي من أجلها لا تتجاوز الكلفة الهامشية $4$ ملايين دينار؟
\qb $(\star)$ نذكّر أن دالة الكلفة الإجمالية $C_T$ هي دالة أصلية لدالة الكلفة الهامشية. أحسب القيمة المتوسطة للكلفة الهامشية عندما تنتج الشركة ما بين $1000$ وحدة و $4000$ وحدة.
\end{sq}
\end{tamrin}

\begin{tamrin}{تمرين من بكالوريا 2017 -- الموضوع الثاني}
نعتبر الدالة $f$ المعرفة على $D_f=]-\infty;0[\cup]0;+\infty[$ كما يلي: $f(x)=\dfrac12e^x-\dfrac1{e^x-1}$ ، و $(C_f)$ تمثيلها البياني.
\QQ{1}
\begin{sq}
\qa أحسب النهايات: $\Lm{-\infty}f(x)$ ، $\Lm{0^-}f(x)$ و $\Lm{0^+}f(x)$ ، وفسّر بيانيًا النتائج المحصل عليها.
\qb أحسب النهاية $\Lm{+\infty}f(x)$.
\end{sq}
\QQ{2}
\begin{sq}
\qa بيّن أنه من أجل كل $x$ من $D_f$: $f'(x)=\dfrac12e^x+\dfrac{e^x}{(e^x-1)^2}$.
\qb استنتج اتجاه تغير الدالة $f$ ثم شكّل جدول تغيراتها.
\end{sq}
\QQ{3} ادرس الوضعية النسبية للمنحنى $(C_f)$ مع المستقيم $(\Delta)$ ذي المعادلة $y=1$.
\QQ{4} عيّن معادلة لـ $(T)$ المماس للمنحنى $(C_f)$ عند النقطة ذات الفاصلة $\ln3$.
\QQ{5} نعتبر الدالة $g$ المعرفة على $]0;+\infty[$ كما يلي: $g(x)=f(x)-\dfrac94(x-\ln3)-1$ ، وهي متزايدة تمامًا من $-\infty$ إلى $+\infty$.
\begin{sq}
\qa أحسب $g(\ln3)$ واستنتج إشارة $g(x)$ حسب قيم $x$.
\qb ادرس على المجال $]0;+\infty[$ وضعية المنحنى $(C_f)$ بالنسبة إلى المماس $(T)$ ، ثم فسّر ذلك بيانيًا.
\end{sq}
\QQ{6} أحسب $f(\ln2)$ ثم ارسم المماس $(T)$ و $(C_f)$ على المجال $]-\infty;0[\cup]0;3]$.
\end{tamrin}

\begin{tamrin}{تمرين من بكالوريا 2017 (الدورة الاستثنائية) -- الموضوع الأول}
\textbf{I)} لتكن $g$ الدالة المعرفة على $\R$ كما يلي: $g(x)=2x-1-e^{2x}$.
\QQ{1} ادرس اتجاه تغير الدالة $g$.
\QQ{2} استنتج إشارة $g(x)$.\par
\textbf{II)} نعتبر الدالة $f$ المعرفة على $\R$ كما يلي: $f(x)=x^2-x-\dfrac12e^{2x}$ ، و $(C_f)$ تمثيلها البياني ($\|\vec i\|=2\,cm$).
\QQ{1} أحسب $\Lm{-\infty}f(x)$ و $\Lm{+\infty}f(x)$. (يعطى: $\Lm{+\infty}\dfrac{e^{2x}}{x^2}=+\infty$).
\QQ{2} ادرس اتجاه تغير الدالة $f$ ثم شكّل جدول تغيراتها.
\QQ{3}
\begin{sq}
\qa بيّن أن المعادلة $f(x)=0$ تقبل حلًا وحيدًا $\alpha$ حيث $-0{,}25<\alpha<-0{,}24$.
\qb أثبت أن المنحنى $(C_f)$ يقبل نقطة انعطاف $A$ إحداثياها $\left(0;-\dfrac12\right)$.
\qc اكتب معادلة المماس $(T)$ للمنحنى $(C_f)$ في النقطة $A$.
\end{sq}
\QQ{4} ارسم $(T)$ و $(C_f)$.
\QQ{5} $(\star)$
\begin{sq}
\qa أحسب بالسنتيمتر المربع المساحة $A(\alpha)$ للحيز المحدد بالمنحنى $(C_f)$ والمستقيمات التي معادلاتها $y=0$ ، $x=\alpha$ و $x=0$.
\qb تحقق أن $A(\alpha)=\dfrac13\left(4\alpha^3-12\alpha^2+6\alpha+3\right)cm^2$.
\end{sq}
\end{tamrin}

% =====================================================
\section{أسئلة الاختيار من متعدد (من البكالوريا والاختبارات)}
\begin{tamrin}{تمارين الاختيار من متعدد}
عيّن الاقتراح الصحيح الوحيد من بين الاقتراحات الثلاثة مع التبرير.
\QQ{1} النهاية $\Lm{-\infty}\left(\dfrac{e^x-1}{e^x+1}\right)$ تساوي:\ABC{$-1$}{$1$}{$+\infty$}
\QQ{2} مجموعة حلول المتراجحة $e^{-x}-e^2\geqslant0$ في $\R$ هي:\ABC{$]-\infty;-2]$}{$[-2;2]$}{$[-2;+\infty[$}
\QQ{3} الدالة المشتقة للدالة $f(x)=(1-x)e^x$ هي:\ABC{$f'(x)=1-xe^x$}{$f'(x)=(1-2x)e^x$}{$f'(x)=-xe^x$}
\QQ{4} $(c_k)$ منحنى الدالة $k(x)=e^{1-x}$ ، والمستقيم $(T):\ y=-x+2$ مماس لـ $(c_k)$ في النقطة $A$ ذات الفاصلة $x_0$ حيث:\ABC{$x_0=1$}{$x_0=0$}{$x_0=-1$}
\QQ{5} مجموعة حلول المعادلة $1-3e^{-x}=0$ في $\R$ هي:\ABC{$\{-\ln3\}$}{$\left\{\dfrac13\right\}$}{$\{\ln3\}$}
\QQ{6} المعادلة $2e^{2x}+3e^x-2=0$ ذات المجهول الحقيقي $x$:\ABC{تقبل حلًا وحيدًا}{تقبل حلين مختلفين}{لا تقبل حلولًا}
\QQ{7} $(\star)$ (بكالوريا 2024) $\displaystyle\int_0^1\left(3x^2+3e^{3x}\right)dx$ يساوي:\ABC{$e^3+1$}{$e^3$}{$e^3-1$}
\QQ{8} (بكالوريا 2023) مجموعة الحلول في $\R$ للمعادلة $e^{2x}+4e^x-5=0$ هي:\ABC{$\{0\}$}{$\{1;0\}$}{$\{-5;0\}$}
\QQ{9} (بكالوريا 2023) $\Lm{+\infty}\left(x+1-e^x\right)$ تساوي:\ABC{$-\infty$}{$+\infty$}{$0$}
\QQ{10} مجموعة حلول المعادلة $1-4e^{-2x}=0$ في $\R$ هي:\ABC{$\{\ln2\}$}{$\{-\ln2\}$}{$\left\{-\dfrac18\right\}$}
\QQ{11} $f(x)=\dfrac{e^x-2}{2e^x+1}$ ، النهاية $\Lm{-\infty}f(x)$ هي:\ABC{$-2$}{$2$}{$\dfrac12$}
\QQ{12} $(\star)$ القيمة المتوسطة للدالة $f(x)=xe^{-x}$ على المجال $[0;1]$ هي:\ABC{$1-\dfrac2e$}{$\dfrac1e$}{$1-\dfrac1e$}
\end{tamrin}

\begin{BOX}{cSec}{}
\centering{\Large\bfseries بالتوفيق في شهادة البكالوريا}
\end{BOX}
\end{document}
